\newpage
\section{One-Page Presentation Outline}

%=========================================================

% SPACE ---   \hspace{0.6em}   ---
% GLOSSARY ---   \gls{gt}   ---
% BOLD ---   \textbf{mathematics}   ---

%=========================================================


\noindent\textbf{Research Topic}: \gls{ml} Approaches for Static Bug Detection

\vspace{0.5em}

\noindent\textbf{Presentation Goal:} Research and expose how \gls{ml} can improve traditional \gls{sca} by reducing false positives and helping developers focus on real software defects.

\vspace{0.5em}

\begin{enumerate}
    \item \textbf{Introduction and Motivation}
    \begin{itemize}
        \item Why software quality and early bug detection are important;
        \item The problem of false positives and the time developers spend checking warnings;
        \item How \gls{ml} could help improve this process.
    \end{itemize}

    \item \textbf{Theoretical Background}
    \begin{itemize}
        \item \gls{sca} detects possible defects without running the program;
        \item Sample traditional tools such as \textit{FindBugs} and their limitations;
        \item Using \gls{sl} to classify warnings as real defects or false positives;
        \item Main evaluation metrics: precision, recall, F1-score and accuracy.
    \end{itemize}

    \item \textbf{Research Papers and Main Findings}
    \begin{itemize}
        \item \textbf{Paper 1 -- Learning to reduce false positives in analytic bug detectors -- Kharkar et al. (2022)}:
        Uses Transformer based models to filter \textit{Infer} \cite{meta_infer} tool warnings, improving precision from 72,7\% to 85,4\% with \textit{GPT-C} \cite{paper1_kharkarlearning2022}.

        \item \textbf{Paper 2 -- Static code analysis alarms filtering reloaded -- Hegedus and Ferenc (2022)}:
        Studies over 224.000 \textit{SonarQube} \cite{sonarqube_docs} warnings and uses \gls{ml} classifiers to filter false positives. \gls{rf} achieves approximately 91\% accuracy \cite{paper2_hegedHus2022static}.

        \item \textbf{Paper 3 -- An empirical assessment of machine learning approaches for triaging reports of static analysis tools -- Yerramreddy et al. (2023)}:
        Compares traditional \gls{ml}, \gls{lstm} and \gls{ggnns}. Shows that code representation and data preparation strongly affect classification results \cite{paper3_yerramreddy2023empirical}.
    \end{itemize}

    \item \textbf{Comparison and Limitations}
    \begin{itemize}
        \item \gls{ml} can reduce false positives and make static analysis more useful;
        \item Different models have different strengths, depending on the code and datasets;
        \item Expose found main limitations as incorrectly filtered real defects, training data quality and poor generalization to new projects.
    \end{itemize}

    \item \textbf{Future Work}
    \begin{itemize}
        \item Improve model accuracy and explainability;
        \item Use larger datasets from different projects and programming languages;
        \item Explore \gls{llms} to analyze warnings, explain possible defects and suggest code fixes;
        \item Test these approaches in real software development environments.
    \end{itemize}

    \item \textbf{Conclusion}
    \begin{itemize}
        \item \gls{ml} can support traditional static analysis, but it cannot completely replace it;
        \item Combining static analyzers with \gls{ml} and \gls{llms} could help developers find and fix defects more efficiently.
    \end{itemize}
\end{enumerate}

\newpage